\section{Belief Loops and Fixed-Point Selection}
\label{sec:belief-loop}

Section~\ref{sec:dynamic} produced two local comparisons. Above \(q_{\mathrm{D1}}\), the D1 off-path map attributes an ambitious deviation to the low type. Above \(q_{\mathrm{IC}}\), the high type's anticipatory incentive constraint fails even under full high-type attribution. Each comparison holds the rest of the assessment fixed. This section solves the whole assessment as a closed configuration.

The issue is not a sequence in which one type exits first and another type reacts later. The object is a fixed-point feasibility test: can any psychological assessment make \(s_A\) sustainably occupied at the top while satisfying sender optimality, Bayes' rule on path, and D1 off path?

The three ingredients are already in place. First, adverse repricing creates reverse selection because \(\Phi_H>\Phi_L\) (Section~\ref{sec:surprise}). Second, withdrawal is individually rational for the high type once the continuation exposure price is large enough. Third, when these two facts are evaluated inside a closed belief hierarchy, the ambitious top has no self-consistent occupancy fixed point. The last step is the state-level annihilation result: for \(q>q^*\), there is no fixed point in which the top type sustainably inhabits \(s_A\).

Thus the theorem below is a simultaneous consistency statement. It asks whether a candidate assessment can close the receiver's interpretation, the sender's second-order belief, and both types' participation inequalities at the same time.

\subsection{Belief hierarchy and psychological assessments}
\label{subsec:belief-hierarchy}

The fixed point is taken over a psychological assessment, not over a strategy profile alone. At a low-scrutiny history \(h_t\), let
\[
\sigma_\theta(h_t)\in\{s_A,s_P\}
\]
denote the type-\(\theta\) sender's signal choice. Let
\[
\rho_A(h_t):=\Pr(\theta_H\mid s_A,h_t)
\]
denote the receiver's first-order posterior after observing the ambitious signal. If \(s_A\) is on path, \(\rho_A\) is pinned down by Bayes' rule. If \(s_A\) is off path, \(\rho_A\) is selected by the refinement.

The sender's payoff depends on what she expects the receiver to believe after \(s_A\). Write
\[
\widehat \rho_{\theta,A}(h_t)
\]
for the type-\(\theta\) sender's second-order belief about the receiver's posterior following \(s_A\). A psychological assessment is a triple
\[
\mathcal A=(\sigma,\rho_A,\widehat \rho_A),
\]
where \(\widehat \rho_A=(\widehat\rho_{H,A},\widehat\rho_{L,A})\). Belief consistency requires
\[
\widehat\rho_{\theta,A}=\rho_A
\qquad\text{for each }\theta\in\{\theta_H,\theta_L\}.
\]
Thus the sender's second-order belief about the receiver's first-order posterior must agree with the posterior induced by the assessment itself.

Crucially, the fixed point is evaluated exactly at this second-order hierarchy layer. To formally embed this within the framework of psychological games \citep{battigalli2009dynamic}, the sender's expected utility must be continuous with respect to the relevant belief hierarchy. Here, the net benefit $B_\theta(\widehat{\rho}_{\theta,A})$ depends continuously on the expected first-order posterior. While the D1 refinement injects a discontinuous jump into the \emph{receiver's} belief map $\rho_A(\sigma)$ when the signal becomes off path, the fixed point is resolved on the composition of these expectations, ensuring that the belief loop strictly and continuously closes at $\widehat\rho_{\theta,A} = \rho_A$.

This is the level at which the fixed point is closed. The exposure term \(\Phi_\theta\) remains true-type indexed. It is not changed by the receiver's posterior. The receiver's posterior changes the current reputational benefit of the signal; it does not relabel the deviator's adverse-repricing exposure.

Let
\[
B_\theta(\rho)
\]
denote the type-\(\theta\) current net benefit from choosing \(s_A\) rather than \(s_P\) when the ambitious signal is interpreted with posterior \(\rho\). This term includes the current reputational return from the receiver's interpretation, net of current signal cost and low-scrutiny epistemic cost. It excludes future high-scrutiny exposure. The fixed \(B_H\) used in Section~\ref{sec:dynamic} is the full high-type-attribution value:
\[
B_H=B_H(1).
\]
Assume
\[
B_H(\rho)\text{ is weakly increasing in }\rho,
\qquad
B_H(1)>0,
\qquad
B_H(0)\le 0.
\]
The first inequality says that a higher high-type posterior makes the ambitious signal more valuable. The last inequality says that an ambitious signal interpreted as low-type evidence carries no high-type premium.

For the low type, assume
\[
B_L(0)<0.
\]
Thus if \(s_A\) is interpreted as low-type evidence, the low type does not receive enough current benefit to justify choosing it.

\subsection{Depreciating exposure and the withdrawal threshold}
\label{subsec:depreciating-exposure}

The dynamic withdrawal motive requires exposure to be persistent but not permanently locked. A sender must be able to reduce future exposure by withdrawing from the ambitious signal. To make this explicit, let
\[
e_t\ge 0
\]
denote the sender's exposed stock at the beginning of period \(t\). Sending \(s_A\) adds one unit of new exposed ambition; choosing \(s_P\) adds no new unit. Existing exposure depreciates at rate \(\eta \in (0,1]\):
\[
e_{t+1}=(1-\eta)e_t+\mathbf{1}\{s_t=s_A\}.
\]
The binary exposure specification in Section~\ref{sec:model} is the special case \(\eta=1\), where withdrawal immediately removes the reduced-form exposure state. The general stock formulation makes clear why withdrawal has dynamic value: it stops new exposure from being created and lets old exposure decay.

Let
\[
\Gamma(q,\eta)>0
\]
denote the discounted marginal continuation price of adding one unit of exposed ambition in the current low-scrutiny state. In the one-period reduced form used in Section~\ref{sec:dynamic},
\[
\Gamma(q,\eta)=\beta q\mu_H.
\]
In the stock version, \(\Gamma(q,\eta)\) is the expected discounted value of the future judgment intensities applied to the marginal unit of exposure:
\[
\Gamma(q,\eta)
=
\mathbb E_q\left[
\sum_{\tau\ge 1}
\beta^\tau (1-\eta)^{\tau-1}\mu_{j_{t+\tau}}
\,\middle|\,j_t=L
\right].
\]
The maintained monotonicity condition is
\[
\Gamma(q,\eta)\text{ is continuous and strictly increasing in }q.
\]
This condition says that a higher probability of entering high scrutiny raises the continuation price of current exposure.

Given posterior \(\rho\), the high type's payoff gain from choosing \(s_A\) rather than \(s_P\) is
\[
\Delta_H(\rho,q)
=
B_H(\rho)-\Gamma(q,\eta)\Phi_H.
\]
Define the high-type withdrawal threshold by
\[
q_{\mathrm{IC}}
=
\inf\{q:\,B_H(1)-\Gamma(q,\eta)\Phi_H<0\}.
\]
In the one-period reduced form, this collapses to the threshold in Section~\ref{sec:dynamic}:
\[
q_{\mathrm{IC}}
=
\frac{B_H(1)}{\beta\mu_H\Phi_H}.
\]

\begin{lemma}[Monotone stopping]
\label{lem:monotone-stopping}
Suppose \(B_H(\rho)\) is weakly increasing in \(\rho\), \(\Phi_H>0\), and \(\Gamma(q,\eta)\) is continuous and strictly increasing in \(q\). Then the high type's stay advantage \(\Delta_H(\rho,q)\) is strictly decreasing in \(q\). In particular, if \(q>q_{\mathrm{IC}}\), then
\[
\Delta_H(\rho,q)<0
\qquad\text{for every }\rho\in[0,1].
\]
Thus the high type's withdrawal region is an upper interval in \(q\).
\end{lemma}

\begin{proof}
For fixed \(\rho\),
\[
\Delta_H(\rho,q)=B_H(\rho)-\Gamma(q,\eta)\Phi_H.
\]
Since \(\Phi_H>0\) and \(\Gamma(q,\eta)\) is strictly increasing in \(q\), \(\Delta_H(\rho,q)\) is strictly decreasing in \(q\). Since \(B_H(\rho)\le B_H(1)\) for every \(\rho\in[0,1]\), the inequality
\[
B_H(1)-\Gamma(q,\eta)\Phi_H<0
\]
implies
\[
B_H(\rho)-\Gamma(q,\eta)\Phi_H<0
\]
for every \(\rho\in[0,1]\). Hence, once \(q>q_{\mathrm{IC}}\), no receiver posterior can make the high type willing to stay on \(s_A\).
\end{proof}

Lemma~\ref{lem:monotone-stopping} is the threshold structure needed for the dynamic interpretation. The stopping region is not fragmented because the only continuation object that changes with \(q\) is the marginal exposure price, and that price moves monotonically.

\subsection{D1 and the off-path belief map}
\label{subsec:off-path-belief-map}

The next object is the receiver's belief after an off-path ambitious signal. In a pooling candidate on \(s_P\), the signal \(s_A\) is not observed on path. The receiver must decide whether an off-path \(s_A\) is more plausibly sent by \(\theta_H\) or by \(\theta_L\).

Let the off-path cost of sending \(s_A\) for type \(\theta\) be
\[
K_\theta(q)=\bar k_\theta+\Gamma(q,\eta)\Phi_\theta,
\]
where \(\bar k_\theta\) is the current off-path cost component and \(\Gamma(q,\eta)\Phi_\theta\) is the continuation exposure cost. Given receiver response \(a\), let \(R(a)\) be the reputational return from that response. The deviation gain is
\[
G_\theta(a;q)=R(a)-K_\theta(q),
\]
and the profitable-deviation response set is
\[
\mathcal D_\theta(q)=\{a:G_\theta(a;q)>0\}.
\]

Because \(\Phi_H>\Phi_L\), the high type's off-path cost rises faster with \(q\). Define
\[
q_{\mathrm{D1}}
=
\inf\{q:\,K_H(q)>K_L(q)\}.
\]
For \(q>q_{\mathrm{D1}}\),
\[
K_H(q)>K_L(q),
\]
and therefore
\[
\mathcal D_H(q)\subsetneq \mathcal D_L(q)
\]
under the strict-response richness condition used in Lemma~\ref{lem:d1-flip}. D1 then assigns the off-path ambitious signal to the low type. Hence the D1-selected posterior is
\[
\rho_A=0.
\]

This D1 conclusion is an off-path counterfactual rule. It is obtained by comparing the response sets \(\mathcal D_H(q)\) and \(\mathcal D_L(q)\); it does not presuppose that either type actually deviates. In particular, low-type withdrawal is not a premise of the D1 classification. The classification follows from the cost ordering induced by \(\Phi_H>\Phi_L\) and \(q>q_{\mathrm{D1}}\).

The on-path and off-path belief map can therefore be summarized as follows. Let
\[
z\in\{0,1\}
\]
denote whether the high type is expected to occupy the ambitious signal:
\[
z=1 \quad\text{means } \theta_H \text{ chooses }s_A,
\]
while
\[
z=0 \quad\text{means } \theta_H \text{ withdraws to }s_P.
\]
For \(q>q_{\mathrm{D1}}\), the receiver's belief about \(s_A\) is
\[
\rho_A(z)=
\begin{cases}
1, & \text{if }z=1\text{ and }s_A\text{ is the separating high-type signal},\\
0, & \text{if }z=0\text{ and }s_A\text{ is off path, selected by D1}.
\end{cases}
\]
The first line is Bayes' rule on the separating path. The second line is the D1 off-path counterfactual map.

\subsection{The belief-loop operator}
\label{subsec:belief-loop-operator}

Given the belief map \(\rho_A(z)\), the high type's participation response is
\[
T_q(z)
=
\mathbf{1}\left\{
B_H(\rho_A(z))-\Gamma(q,\eta)\Phi_H\ge 0
\right\}.
\]
A high-type belief-loop fixed point is a value \(z\in\{0,1\}\) such that
\[
z=T_q(z).
\]
This is the self-referential equation. High-type participation determines whether \(s_A\) is on path. The path status of \(s_A\) determines the receiver's posterior. The receiver's posterior determines the current rent of \(s_A\). That rent determines high-type participation.

Define
\[
q^*=\max\{q_{\mathrm{IC}},q_{\mathrm{D1}}\}.
\]

The closed-loop candidate that will survive is the safe-pooling configuration
\[
\mathcal C_0
=
\left(
\sigma_H=s_P,\,
\sigma_L=s_P;\,
s_A\text{ off path};\,
\rho_A^{D1}=0
\right).
\]
Its fixed-point structure is explicit:
\begin{enumerate}
\item Both types choose the protected signal \(s_P\), so \(s_A\) is off path.
\item Because \(s_A\) is off path, D1 evaluates the counterfactual deviation. Since \(\Phi_H>\Phi_L\) and \(q>q_{\mathrm{D1}}\),
\[
\mathcal D_H(q)\subsetneq \mathcal D_L(q),
\]
so D1 selects the low-type attribution \(\rho_A^{D1}=0\).
\item Under \(\rho_A^{D1}=0\), the deviation margins are
\[
\Delta_H^{D1}(q)=B_H(0)-\Gamma(q,\eta)\Phi_H<0,
\qquad
\Delta_L^{D1}(q)=B_L(0)-\Gamma(q,\eta)\Phi_L<0.
\]
The first inequality follows from \(B_H(0)\le 0\) and \(\Gamma(q,\eta)\Phi_H>0\). The second follows from \(B_L(0)<0\) and \(\Gamma(q,\eta)\Phi_L>0\).
\item Therefore no type wants to move from \(s_P\) to \(s_A\). The configuration \(\mathcal C_0\) closes on itself; no chronological claim about who deviates first is used.
\end{enumerate}

\begin{theorem}[Belief-loop fixed-point selection]
\label{thm:belief-loop-fixed-point}
Suppose:
\begin{enumerate}
\item \(B_H(\rho)\) is weakly increasing in \(\rho\), with \(B_H(1)>0\) and \(B_H(0)\le 0\);
\item \(B_L(0)<0\);
\item \(\Phi_H>\Phi_L\), as derived from primitive self-accuracy levels and Bayesian-surprise repricing;
\item \(\Gamma(q,\eta)\) is continuous and strictly increasing in \(q\), with \(\eta \in (0,1]\);
\item \(q>q^*=\max\{q_{\mathrm{IC}},q_{\mathrm{D1}}\}\);
\item off-path beliefs satisfy D1;
\item the analysis is restricted to the binary-signal pure-strategy class.
\end{enumerate}
Then no D1-consistent psychological assessment in this class has stable top-end occupancy of \(s_A\). The unique D1-consistent pure closed-loop assessment is \(\mathcal C_0\):
\[
\theta_H\mapsto s_P,
\qquad
\theta_L\mapsto s_P.
\]
Equivalently, the unique high-type occupancy fixed point is
\[
z^*=0.
\]
\end{theorem}

\begin{proof}
We evaluate pure closed-loop configurations.

First, no configuration with top-end occupancy of \(s_A\) is feasible. If only the high type occupies \(s_A\), then \(s_A\) is the separating high-type signal, so Bayes' rule gives
\[
\rho_A(1)=1.
\]
The high type's payoff gain from choosing \(s_A\) rather than \(s_P\) is
\[
\Delta_H(1,q)=B_H(1)-\Gamma(q,\eta)\Phi_H.
\]
Since \(q>q_{\mathrm{IC}}\), this expression is strictly negative. Hence
\[
T_q(1)=0.
\]
Therefore \(z=1\) is not a fixed point.

If both types choose \(s_A\), then \(s_A\) is pooling rather than separating. Let \(\pi_0\in(0,1)\) be the prior probability of \(\theta_H\). Bayes' rule gives \(\rho_A=\pi_0\). Since \(B_H(\rho)\) is weakly increasing and \(\pi_0<1\),
\[
B_H(\pi_0)\le B_H(1).
\]
Because \(q>q_{\mathrm{IC}}\),
\[
B_H(1)-\Gamma(q,\eta)\Phi_H<0.
\]
Therefore
\[
B_H(\pi_0)-\Gamma(q,\eta)\Phi_H<0.
\]
The high type would prefer \(s_P\), so pooling on \(s_A\) is not a fixed point either. Thus every configuration with top-end occupancy of \(s_A\) is infeasible.

Second, no low-only occupancy configuration is feasible. If only the low type chooses \(s_A\), then \(s_A\) is on path as low-type evidence, so \(\rho_A=0\). The low type's gain is
\[
B_L(0)-\Gamma(q,\eta)\Phi_L<0,
\]
because \(B_L(0)<0\) and \(\Gamma(q,\eta)\Phi_L>0\). Hence the low type does not wish to occupy \(s_A\) under low-type attribution.

Third, evaluate \(\mathcal C_0\). Under \(\mathcal C_0\), \(s_A\) is off path. Since \(q>q_{\mathrm{D1}}\), the D1 response-set comparison gives
\[
\mathcal D_H(q)\subsetneq \mathcal D_L(q),
\]
so the off-path ambitious signal is assigned to the low type and
\[
\rho_A^{D1}=0.
\]
The high-type deviation margin is
\[
\Delta_H^{D1}(q)=B_H(0)-\Gamma(q,\eta)\Phi_H<0,
\]
and the low-type deviation margin is
\[
\Delta_L^{D1}(q)=B_L(0)-\Gamma(q,\eta)\Phi_L<0.
\]
Thus neither type has a profitable deviation to \(s_A\). The safe-pooling configuration is internally consistent. Since all other pure configurations have been ruled out, \(\mathcal C_0\) is the unique D1-consistent psychological assessment in the binary pure-strategy class.
\end{proof}

\subsection{The closed loop as a state condition}
\label{subsec:closed-loop}

The theorem should be read as a state claim:
\[
\mathcal C_0
\Longleftrightarrow
\left[
s_A\text{ off path}
\right]
\wedge
\left[
\rho_A^{D1}=0
\right]
\wedge
\left[
\Delta_H^{D1}(q)<0
\right]
\wedge
\left[
\Delta_L^{D1}(q)<0
\right].
\]
The D1 component is not caused by an observed low-type deviation. It is the refinement's counterfactual response to an off-path signal, determined by \(\mathcal D_H(q)\subsetneq\mathcal D_L(q)\). The participation inequalities then show that the counterfactual off-path interpretation is consistent with nobody actually sending \(s_A\).

This separation removes the circular timing story. Reverse selection is the cost-ranking fact \(\Phi_H>\Phi_L\). Individually rational withdrawal is the high-type inequality \(\Delta_H(1,q)<0\) once \(q>q_{\mathrm{IC}}\). Annihilation is the fixed-point conclusion that no D1-consistent pure assessment contains stable top-end occupancy of \(s_A\) once \(q>q^*\).

In the coordination and self-fulfilling-belief literature, forward-induction refinement characteristically selects \emph{toward} efficiency: it eliminates the off-path beliefs that sustain coordination failure, uniquely selecting the Pareto-dominant equilibrium \citep[e.g., the refinement resolutions of][]{diamond1983bank}. This standard selection presumes that the party able to orchestrate a deviation is the one seeking the efficient outcome, so forward induction attributes the deviation's intent accordingly. We exhibit the opposite operation. When malicious scrutiny reverses the cost ranking so that the wider profitable-deviation set belongs to the \emph{low} type, the same D1 logic confirms rather than dissolves the absence of stable top-end occupancy: the refinement selects safe pooling. This resulting top-down pooling is an \emph{informational} disaster, destroying the market's ability to screen types. Yet it is also privately rational risk avoidance for high types facing excessive epistemic exposure. The refinement does not rescue the informational efficiency of the separating equilibrium; instead, it enforces the nonexistence of an occupied ambitious-top fixed point.

\subsection{No interior mixing in this section}
\label{subsec:no-mixing-sec6}

The fixed-point theorem is a pure-strategy selection result. It does not use an interior mixing probability for the high type. This restriction is deliberate. In a binary model with no low-type background arrival at \(s_A\), any strictly positive high-type mass on \(s_A\) would keep the posterior after \(s_A\) degenerate in the high type. Continuous signal dilution therefore cannot be generated by high-type mixing alone.

Interior mixing requires an additional object, such as low-type mimetic entry, noisy audience inference, trembles, or a continuous type distribution. We explicitly choose to scope out mixing in this section to isolate the exact conditions under which state-level collapse of top-end occupancy occurs in the binary pure-strategy class. The interior mixed regime reemerges naturally as the \emph{Polarization} phenomenon handled in Section~\ref{sec:polarization} (and the continuous-type phase transition), where the signal informativeness continuously degenerates without dying entirely. The present section proves only the sharper fixed-point claim: above \(q^*\), the only D1-consistent pure belief loop is \(\mathcal C_0\), so the ambitious signal has no stable top-end occupancy fixed point.
